% resumen.tex: BORRADOR (2026-10-18, revisado el 2026-10-19) para revisión del autor.
% fuente: H-031 (concentración), H-039/H-040 (separabilidad), H-042 (AUC del MLP), H-044 (QSVM),
%         H-048 (etiquetado diseñado; X6_etiquetas_construidas.csv), H-052 (OE-6). Todas las cifras están también en §6.
En este reporte se evalúa la viabilidad de las transformaciones cuánticas variacionales como mecanismo de codificación de características radiómicas extraídas de mamografías reales. Sobre las 3,568 lesiones del CBIS-DDSM, con masas y calcificaciones como subproblemas independientes, se extraen 67 características con PyRadiomics, se seleccionan 12 y se escalan a $[0,\pi]$. El mismo vector alimenta cinco condiciones: sin transformación, PCA, Kernel PCA con kernel RBF y dos \textit{feature maps} cuánticos (ZZ y Pauli con términos Y y ZZ) seguidos de un \textit{ansatz} de parámetros fijos, simulados con Qiskit Aer. La separabilidad se mide sobre los kernels y sobre los vectores (\textit{kernel-target alignment}, \textit{geometric difference} regularizada, dimensión efectiva, Davies-Bouldin y razón de Fisher), y la clasificación con un mismo perceptrón multicapa y validación cruzada por paciente.

Con las etiquetas diagnósticas, las representaciones cuánticas no muestran mayor separabilidad en ninguna métrica. El kernel cuántico está concentrado (mediana fuera de la diagonal cercana a 0.003, frente a 0.61 del RBF) y la \textit{geometric difference} (1.40 a 1.70) queda muy por debajo del umbral $\sqrt{N}\approx 14$. El perceptrón pierde entre 0.07 y 0.14 de AUC-ROC sobre los \textit{embeddings} cuánticos, mientras que una SVM sobre el kernel cuántico iguala a la clásica, lo que sitúa la pérdida en la lectura de los valores esperados y no en la codificación. La ausencia de ventaja depende del problema y no del protocolo: con un etiquetado diseñado según el procedimiento de Huang et al., que alinea las clases con la geometría del kernel cuántico, el mismo protocolo sí detecta la ventaja: el kernel cuántico proyectado alcanza un AUC de 0.998 en masas y 0.994 en calcificaciones, frente a 0.800 y 0.824 del mejor clasificador clásico. Con 1,024 disparos, entre el 16 y el 74\,\% de los valores del kernel se leerían como cero.

\noindent\textbf{Palabras clave:} aprendizaje automático cuántico, kernels cuánticos, separabilidad de clases, radiómica, CBIS-DDSM, cáncer de mama.
